% !TEX root = ../Main.tex
\chapter{分裂法证明——正整数 $\alpha$ 次方的求和}

要求 $1^\alpha+2^\alpha+\cdots+n^\alpha$。$\alpha=1$ 用上一章的倒序就够了；$\alpha\ge 2$ 则用\textbf{分裂法（裂项）}：把 $k^\alpha-(k-1)^\alpha$ 拆开，叠加后让低一次的和自己冒出来。

\section{$\alpha=1$}

$S=1+2+\cdots+n$，同前一章分裂（倒序相加）：$S=n+(n-1)+\cdots+1$，对应相加 $2S=(n+1)\cdot n$，故
\[
S=\frac{n(n+1)}{2}.
\]

\section{$\alpha=2$}

要求 $T=1^2+2^2+3^2+\cdots+n^2$。我们这样想：考察相邻平方的差
\[
\begin{aligned}
n^2-(n-1)^2&=n^2-(n^2+1-2n)=2n-1,\\
(n-1)^2-(n-2)^2&=2n-3=2(n-1)-1,\\
&\ \ \vdots\\
2^2-1^2&=3=2(2)-1,\\
1^2-0^2&=1=2(1)-1.
\end{aligned}
\]
这 $n$ 个式子，\textbf{左右分别求和}。左边逐项抵消（裂项相消）只剩 $n^2-0^2=n^2$，右边是 $2(1+2+\cdots+n)-n$：
\[
n^2=2\,(1+2+\cdots+n)-n
\quad\Longrightarrow\quad
1+2+\cdots+n=\frac{n^2+n}{2}.
\]
注意：从\textbf{二次}的差出发求和，结果落在了\textbf{一次}的求和上——左边比右边高一次，求和时高次大量相消，反倒把低一次的和"逼"了出来，而且式子更简洁。

\section{$\alpha=2$ 的真正求法：升一阶}

上一步的教训是：要得到\textbf{二次和} $T$，就该从\textbf{三次}的差下手。照搬刚才的写法：
\[
\begin{aligned}
n^3-(n-1)^3&=3n^2-3n+1,\\
(n-1)^3-(n-2)^3&=3(n-1)^2-3(n-1)+1,\\
&\ \ \vdots\\
2^3-1^3&=7=3(2)^2-3(2)+1,\\
1^3-0^3&=1=3(1)^2-3(1)+1.
\end{aligned}
\]
左右分别相加。左边裂项相消得 $n^3$；右边是 $3T-3S+n$，其中 $S=1+2+\cdots+n=\dfrac{n(n+1)}{2}$：
\[
n^3=3T-3S+n.
\]
代入 $S$ 解出 $T$：
\[
3T=n^3-n+\frac{3n(n+1)}{2}=\frac{2n^3+3n^2+n}{2}=\frac{n(n+1)(2n+1)}{2},
\]
故
\[
\boxed{\,1^2+2^2+\cdots+n^2=\frac{n(n+1)(2n+1)}{6}\,}.
\]

\section{类比到 $\alpha=3$}

同样的招数升到\textbf{四次}差：取 $(k+1)^4-k^4=f(k)$ 是关于 $k$ 的三次式，左右分别求和，左边裂项相消、右边出现 $\sum k^3$，再代入已知的 $\sum k^2,\sum k$，即可解出
\[
1^3+2^3+\cdots+n^3=\br{\frac{n(n+1)}{2}}^2.
\]
正好是 $1+2+\cdots+n$ 的平方——立方和等于"和的平方"，形式格外漂亮。
